\ficha{Germer (s.d.), A gestação do padrão ouro-câmbio antes de 1914}{germerouro}

\begin{identificacao}
GERMER, Claus Magno. \textbf{A gestação dos elementos do `padrão ouro-câmbio'
antes de I Guerra Mundial e sua institucionalização no pós-guerra}. [S.l.:
s.n.], [s.d.]. 22 p.

Artigo submetido a painel de evento acadêmico, com o resumo de inscrição anexado
ao final (p. 20 a 22), assinado no Departamento de Economia da UFPR. O PDF não
traz veículo, data nem número de edição. A bibliografia cita Fischer (1999) e a
nota 2 remete às crises financeiras do segundo semestre de 1998 (p. 2), o que
situa a redação em 1999 ou depois. Lido o texto inteiro, 22 páginas, em
português, na cópia digital do acervo.
\end{identificacao}

\bloco{Problema e tese}
O texto pergunta quando e por que se formou o padrão ouro-câmbio. A resposta
recusa a periodização corrente, que faz da I Guerra Mundial o divisor de águas.
Os elementos constitutivos do sistema, reservas em moeda estrangeira conversível
e liquidação de saldos por dinheiro de crédito, já operavam sob o chamado
padrão-ouro internacional. A Conferência de Gênova, em 1922, não improvisou um
padrão novo, reconheceu e institucionalizou uma tendência anterior. A causa não
está na guerra, e sim no desenvolvimento dos sistemas de crédito nacionais
integrados, coroados pelo banco central, e na formação de vínculos creditícios
entre eles.

\bloco{Argumento}
\begin{enumerate}[nosep]
\item A periodização corrente usa como critério a forma do meio circulante,
porque se apoia na teoria quantitativa. O modelo de Hume supõe circulação
metálica pura e exclui o crédito, de modo que os autores não separam elementos
monetários e creditícios (p. 1). O elemento novo do fim do século 19 não é o
ouro, que já cumpria funções monetárias havia muito e fazia da Inglaterra um
país objetivamente em padrão-ouro desde 1717 (p. 3), e sim a constituição de
sistemas bancários nacionais integrados dotados de banco central (p. 4).
\item O sistema ouro-câmbio é definido pelo lado das reservas: ao lado do ouro,
os países comuns usam como meio de reserva e de pagamento internacional títulos
denominados no padrão monetário dos centros de emissão, ou centros-ouro (p. 2).
Isso exige uma estrutura de tipo bancário no plano internacional.
\item A precedência é sustentada por autoridade, não por dados próprios. Keynes
já em 1913 usava a expressão e apontava a tendência como dominante (p. 4). Wood
fala de sistema ouro-câmbio libra, e Eichengreen acrescenta duas distinções:
antes da guerra o sistema existia mas não estava institucionalizado, e
funcionava com múltiplas moedas de reserva (p. 4). Gênova recomendou criar
bancos centrais onde ainda não houvesse (p. 4).
\item A explicação das reservas em divisas é o centro do texto. O viés
quantitativista as atribui à necessidade de proteger a reserva de ouro, mas
Germer objeta que o padrão-ouro, na lógica de Hume, não exige reserva
centralizada nem centro financeiro (p. 11). A formação de reservas em moedas
conversíveis reflete duas coisas distintas: a integração dos sistemas bancários
nacionais, cuja reserva metálica é o fundamento da circulação do dinheiro de
crédito, e a formação de um sistema de crédito internacional (p. 12). Os países
secundários passaram a manter carteira de títulos em libra e em dólar para
cobrir déficits momentâneos, evitando ao mesmo tempo a alta do juro interno e a
drenagem de ouro (p. 13).
\item O sistema de crédito tem forma de pirâmide. Em cada andar os saldos
devedores são cobertos com títulos do andar imediatamente superior, e só os
saldos finais são pagos em ouro, nas compensações interbancárias do banco
central (p. 14). Daí a centralização da reserva e a proibição legal dos
vazamentos, processo concluído nos anos 1930.
\end{enumerate}

\chave{Consequentemente, não fazem distinção adequada entre os elementos puramente monetários e os creditícios presentes no sistema}{p. 1}

\bloco{Conceitos e definições operacionais}
Dinheiro é a mercadoria que funciona como equivalente geral, medida de valor e
base do padrão de preços, funções que o ouro não perde ao deixar de circular.
Dinheiro de crédito é o conjunto das formas que circulam em nome do dinheiro, da
letra de câmbio à nota bancária e ao depósito. Sistema de crédito, expressão que
Germer toma de Marx, substitui sistema monetário como designação da esfera da
circulação. Padrão ouro-câmbio é o arranjo em que a reserva combina ouro e
títulos em moeda conversível dos centros de emissão. Padrão-ouro internacional
deixa de ser fase terminal de um sistema metálico e passa a ser a fase inicial do
sistema de crédito internacional, com base monetária em ouro e meio circulante
creditício, ainda segmentado em sistemas nacionais (p. 11). O texto não define
currency board, caixa de conversão, paridade, credibilidade nem regra monetária.

\bloco{Evidência e método}
Ensaio teórico e histórico, sem série, sem tabela e sem contrafactual. A
evidência é literatura secundária e exegese de Marx, com apoio recorrente em
Rist, De Cecco, Eichengreen, Bagehot e Keynes. Os poucos casos concretos são
ilustrativos: a safra agrícola norte-americana drenando ouro para as mãos dos
agricultores por causa da fragmentação bancária, e os bancos de depósito ingleses
que formaram reserva de ouro autônoma frente ao Banco da Inglaterra (p. 15); os
bancos privados franceses que mantinham títulos ingleses em carteira, sobre os
quais Rist observa que Londres, e não Paris, lhes servia de banco central
(p. 13).

\chave{a utilização das divisas como reservas de disponibilidades internacionais era prática corrente antes de 1914}{Niveau, citado na p. 4}

\bloco{Agentes e vertentes}
Dois campos teóricos. De um lado a teoria quantitativa, de Hume a seus
sucessores, com a qual o autor identifica a historiografia corrente do
padrão-ouro. De outro a banking school e Marx, que distinguem dinheiro e dinheiro
de crédito, com Rist como elo. Os pós-keynesianos, Wray e Moore, ficam no erro
simétrico ao quantitativista, por reduzirem o dinheiro ao dinheiro de crédito.
Keynes aparece como precursor da tese, em 1913 e no Treatise. Eichengreen, De
Cecco, Bordo, Wood, Niveau e Goodhart entram como testemunhas do fenômeno, não
como adversários. O autor é professor da UFPR.

\bloco{Críticas e limites}
O texto não trata do Brasil e não menciona caixas de conversão, currency boards
nem moedas coloniais em nenhum ponto. A única ligação com o mundo colonial é
indireta e bibliográfica, a citação de \textit{Indian currency and finance} como
origem da expressão padrão ouro-câmbio (p. 4). A periferia aparece uma vez, de
passagem, na afirmação de que a integração bancária avançava aceleradamente nos
demais países capitalistas, mesmo na periferia do sistema (p. 10), sem exemplo e
sem fonte. O autor não distingue a reserva em divisas de país credor da de país
devedor, que é justamente a distinção de que o caso brasileiro precisa. A
explicação é funcionalista: centralização de reservas e criação de bancos
centrais aparecem como resultado necessário de uma tendência, sem exame das
disputas políticas que as produziram. Não há mensuração da composição das
reservas, de modo que a extensão da prática antes de 1914 fica sustentada apenas
por citação de terceiros.

\begin{dialogo}
\item O que a lei e o debate admitiam como lastro da Caixa de Conversão.
Procurar nos jornais e nos Documentos Parlamentares os termos ``letras'',
``cambiais'', ``saques sobre Londres'' e ``depósito'', ao lado de ``ouro em
barra'' e ``moeda estrangeira''. Se o debate tratava título em libra como
equivalente do metal, o caso brasileiro passa a ser instância periférica da
prática que Germer atribui aos países secundários (p. 12 e 13), lado que ele não
examina.
\item A custódia do ouro em 1911 e a agência de Londres. Separar quem exigia o
metal no Rio de quem aceitava saldo em Londres. A hipótese a testar é que o
argumento de época, ouro em Londres é crédito e não metal, reproduz no
vocabulário dos jornais a distinção entre elemento monetário e elemento
creditício que Germer diz ausente na teoria quantitativa (p. 1).
\item O vocabulário técnico disponível entre 1906 e 1914. Verificar se
articulistas e parlamentares empregavam algo como padrão ouro-câmbio, ou se
falavam apenas em conversibilidade, lastro e taxa fixa. O achado testa a
afirmação de que a expressão de Keynes circulava desde 1913 e de que a prática
precedia o nome (p. 4).
\item A suspensão de 1914. Observar se a defesa da Caixa apelou a saldos e
créditos no exterior. Se apelou, o desenho brasileiro dependia de vínculo
creditício internacional e não só de metal em cofre, que é a tese do texto.
\end{dialogo}

\usonocapitulo{Parte I, como apoio. Serve à afirmação de que o padrão-ouro
anterior a 1914 já operava com reservas em divisas e liquidação por crédito, o
que retira da Caixa de Conversão o estatuto de anomalia periférica e permite
lê-la como variante de um arranjo em formação. O texto não fala do Brasil, não
descreve nenhuma caixa de conversão e não trata de moedas coloniais, então a
comparação institucional tem de vir de Hanke e Schuler e de della Paolera e
Taylor. Daqui vem o argumento conceitual, e sobretudo a distinção entre elemento
monetário e elemento creditício, útil para ler o debate sobre a custódia do ouro
e sobre a agência de Londres.}
